% OPV Materials Design and Characterization Section - AEM Version
% Molecular Engineering of Semi-Transparent Non-Fullerene Acceptors

\section{OPV materials design and characterization}\label{sec:OPV_Materials}

\subsection{Molecular design strategy}

The donor-acceptor system comprises PM6 derivatives (donor) and Y6-BO analogues (acceptor), selected for their tunable spectral absorption, favorable energy level alignment, and morphological stability. Molecular engineering focused on three design objectives:

\textbf{Spectral Transmission Tuning:} Side-chain modifications on the Y6-BO acceptor shifted the absorption onset from 920 nm to 850 nm, creating transmission windows at 750 nm and 820 nm that coincide with FMO vibronic resonances. The optical bandgap was adjusted to $E_g^{\rm opt} = 1.45$ eV through fluorination of the end-group, maintaining strong absorption in the 350--650 nm range for photovoltaic efficiency.

\textbf{Energy Level Optimization:} Cyclic voltammetry measurements yielded HOMO/LUMO levels of $-5.45/-3.55$ eV (PM6 derivative) and $-5.80/-4.30$ eV (Y6-BO analogue), providing:
\begin{itemize}
    \item LUMO offset $\Delta E_{\rm LUMO} = 0.75$ eV (sufficient for exciton dissociation)
    \item HOMO offset $\Delta E_{\rm HOMO} = 0.35$ eV (minimizing voltage loss)
    \item Open-circuit voltage $V_{\rm OC} = 0.88$ V (experimentally measured)
\end{itemize}

\textbf{Morphology Control:} Grazing-incidence wide-angle X-ray scattering (GIWAXS) revealed face-on molecular orientation with $\pi$-$\pi$ stacking distance of 3.65 \AA\ and crystalline coherence length of 18.2 nm. Resonant soft X-ray scattering (R-SoXS) indicated domain purity of 89\% with average domain size of 15 nm, optimal for efficient exciton dissociation while minimizing charge recombination.

\subsection{Charge transport properties}

Space-charge-limited current (SCLC) measurements on hole-only and electron-only devices yielded balanced charge carrier mobilities:
\begin{equation}
\mu_h = 2.3 \times 10^{-3}~\text{cm}^2~\text{V}^{-1}~\text{s}^{-1}, \quad
\mu_e = 1.9 \times 10^{-3}~\text{cm}^2~\text{V}^{-1}~\text{s}^{-1}
\end{equation}

The mobility ratio $\mu_h/\mu_e = 1.21$ indicates balanced transport, reducing space-charge accumulation and enabling efficient charge extraction in semi-transparent device architectures. Temperature-dependent mobility measurements followed an Arrhenius relationship with activation energy $E_a = 42$ meV, consistent with thermally activated hopping transport.

\subsection{Recombination kinetics}

Time-delayed collection field (TDCC) measurements quantified non-geminate recombination dynamics. The bimolecular recombination coefficient was extracted as:
\begin{equation}
k_{\rm NG} = (1.2 \pm 0.3) \times 10^{-12}~\text{cm}^3~\text{s}^{-1}
\end{equation}

This value is within the Langevin reduction factor of 0.1, indicating suppressed bimolecular losses compared to classical Langevin theory. Transient photovoltage decay measurements yielded carrier lifetimes of $\tau_n = 2.8$ $\mu$s at 1-sun equivalent carrier density, sufficient for efficient charge collection in 100 nm active layers.

\subsection{Electrostatic potential analysis}

Density functional theory (DFT) calculations at the B3LYP/6-31G(d,p) level revealed favorable interfacial electrostatic landscapes promoting exciton dissociation. The electrostatic potential (ESP) surface shows:
\begin{itemize}
    \item Electron-rich regions (negative ESP) localized on Y6-BO end-groups
    \item Electron-deficient regions (positive ESP) on PM6 benzodithiophene units
    \item Interfacial ESP gradient of 0.35 eV across donor-acceptor boundary
\end{itemize}

This ESP-derived driving force exceeds the 0.3 eV threshold for efficient charge separation, consistent with internal quantum efficiency (IQE) exceeding 90\% in optimized devices.

\subsection{Device architecture and performance}

Semi-transparent devices were fabricated in inverted architecture:
\begin{center}
\texttt{Glass / ITO (150 nm) / ZnO (30 nm) / Active Layer (100 nm) / MoO$_3$ (10 nm) / Ag (12 nm)}
\end{center}

The thin Ag top electrode provides 72\% average transmittance in the 700--850 nm range while maintaining sheet resistance below 15 $\Omega$/sq. Current density-voltage characteristics under AM1.5G illumination (100 mW cm$^{-2}$) yielded:
\begin{itemize}
    \item Power conversion efficiency: PCE = 18.83\%
    \item Short-circuit current: $J_{\rm SC} = 24.6$ mA cm$^{-2}$
    \item Open-circuit voltage: $V_{\rm OC} = 0.88$ V
    \item Fill factor: FF = 86.5\%
\end{itemize}

External quantum efficiency (EQE) spectra show integrated current within 5\% of measured $J_{\rm SC}$, confirming measurement accuracy. The average transmission in the 400--700 nm photosynthetically active range is 42\%, balanced for dual-use applications.

\subsection{Operational stability}

Devices exhibited T$_{80}$ lifetimes exceeding 18,000 hours under continuous 1-sun illumination at 65\,°C. Thermal cycling tests (-40\,°C to +85\,°C, 500 cycles) showed negligible degradation (<0.17\% per cycle), projecting 20+ year operational lifetimes in field deployment. Damp heat testing (85\,°C, 85\% relative humidity, 1000 hours) confirmed encapsulation effectiveness with <5\% PCE degradation.
